% Generated draft: requires mathematical and source review.
\documentclass[11pt]{article}
\usepackage[margin=25mm]{geometry}
\usepackage{fontspec}
\setmainfont{Latin Modern Roman}
\usepackage{amsmath,amssymb,mathtools,unicode-math}
\setmathfont{Latin Modern Math}
\usepackage{enumitem,needspace,xurl,hyperref}
\setlength{\parindent}{0pt}
\setlength{\parskip}{5pt}
\setlength{\emergencystretch}{3em}
\newcommand{\N}{\mathbb{N}}
\newcommand{\Z}{\mathbb{Z}}
\newcommand{\Q}{\mathbb{Q}}
\newcommand{\R}{\mathbb{R}}
\newcommand{\C}{\mathbb{C}}
\newcommand{\F}{\mathbb{F}}
\newcommand{\A}{\mathbb{A}}
\newcommand{\PP}{\mathbb{P}}
\newcommand{\E}{\mathbb{E}}
\newcommand{\eps}{\varepsilon}
\newcommand{\dd}{\,\mathrm{d}}
\begin{document}
\section*{UnsolvedMath ID: 20003279}
\textbf{Source problem code:} catalogue entry
\subsection{Definitions and mathematical statement}
Question 12, attributed to F. Ledrappier and O. Sarig, uses the source word ``conilpotent.'' The aid reports that the official PDF really prints this word and interprets it as follows: $G=SL(2,\R)$, $\Lambda\subset G$ is a uniform lattice, $\Gamma\triangleleft\Lambda$, and the deck group $\Lambda/\Gamma$ is nilpotent. This is a quotient condition, not nilpotence of $\Gamma$ or an unstated residual-nilpotence condition. Let $Y=G/\Gamma$, let $m$ be quotient Haar measure, and let $U=\{u_t\}$ be the left horocycle flow with $u_t$ upper triangular unipotent.

In the full-flow centralizer convention used by the aid, squashability means existence of an invertible nonsingular measurable $S:Y\to Y$ commuting with every $u_t$, modulo null sets, such that $S_*m=c\,m$ for a constant $c>0$, $c\ne1$. Nonsquashability means no such map exists. The question asks whether $m$ is nonsquashable for all the stated nilpotent regular covers. The aid reports the coabelian case known and the finite-index case positive, while finding no resolution of the infinite nonabelian nilpotent case through August 14, 2026. It does not extend that search conclusion to an independently checked present-day status.

The aid's general reduction assumes an ergodic measure-preserving $U$ action on a sigma-finite measure space. For any invertible nonsingular measurable map commuting with the full flow, the Radon--Nikodym derivative of $S_*m$ is $U$-invariant and hence a positive constant. Thus
\[
 \chi:\operatorname{Cent}_{\mathrm{ns}}(U)\longrightarrow\R_{>0},
 \qquad S_*m=\chi(S)m,
 \qquad\chi(ST)=\chi(S)\chi(T)
\]
is a scaling character. Nonsquashability is triviality of this character. Ergodicity and full-flow commutation are hypotheses; a time-one centralizer or a map merely normalizing time is a different object.

For quotient Haar measure, all explicitly quotient-defined examples in the aid have multiplier one: left translations by elements of $C_G(U)$, right translations by $N_G(\Gamma)$, and in particular deck maps from $\Lambda/\Gamma$. The right maps are defined by $g\Gamma\mapsto gn\Gamma$ for $n$ normalizing $\Gamma$, and preserve Haar by unimodularity and the quotient normalization. Therefore a nontrivial value of $\chi$, when the ergodic reduction applies, must come from a measurable centralizer map beyond these examples. If $m(Y)<\infty$, invertibility already forces any constant multiplier to equal one.

Equivalently, form the graph measure $\eta_S=(\mathrm{id},S)_*m$ on $Y\times Y$. Commutation makes it diagonally $U$-invariant, and
\[
 (p_1)_*\eta_S=m,\qquad(p_2)_*\eta_S=\chi(S)m.
\]
A squash gives unequal proportional Haar marginals. Conversely a diagonally invariant graph quasi-joining whose graph defines an invertible nonsingular map and whose two marginals are unequal positive multiples of the fixed Haar measure gives such a squash after normalization. Invertibility and graph support are essential; an arbitrary joining with prescribed marginals is not the claimed equivalence. The required Haar ergodicity and standard measurable conventions must be matched to the cover setting.

The following source prose is separate context: it raises analogous historical questions for general discrete $\Gamma$ and other $U$-ergodic Radon measures, then turns to higher-rank abelian measure rigidity. It describes positive-entropy hypotheses in [121] and attributes to E. Lindenstrauss [122] uniqueness of Haar probability among $D(\mathbb A)$-invariant measures on $SL(2,\mathbb A)/SL(2,\Q)$, where $\mathbb A$ is the adele restricted product and $D$ the diagonal group. This adjacent adelic assertion is not an added hypothesis or conclusion of the nilpotent-cover question.
\subsection{Short English statement}
For a regular nilpotent cover of a compact hyperbolic quotient, must every measurable horocycle-centralizing scaling map preserve Haar measure? The supplied aid reduces this to triviality of a centralizer character, equivalently absence of invertible graph quasi-joinings with unequal Haar marginals. It rules out the explicit translation and deck maps as squashes, while leaving the infinite nonabelian nilpotent case unresolved in its dated check.
\subsection{Sources}
\begin{itemize}
\item Local UnsolvedMath registry, permanent ID $20003279$, source code catalogue entry; canonical location \texttt{aim-other-notes.json}, \texttt{notes[85]}.
\item AIM \emph{Emerging applications of measure rigidity}, Question 12 (F. Ledrappier, O. Sarig), \url{https://aimath.org/WWN/measrigid/measrigid.pdf}. The supplied verification says ``conilpotent'' is printed in the PDF, not merely OCR.
\item Ledrappier--Sarig comparison and coabelian results are mentioned in the attached aid without full references. Its literature check through August 14, 2026 found no solution of the infinite nonabelian nilpotent case.
\item The appended section cites [121] for higher-rank abelian measure rigidity and [122] for E. Lindenstrauss's stated adelic result. Full bibliography entries are not supplied.
\item Research attempt 1 is marked reduction, proved, and low novelty confidence, with an explicit machine-generated, non-peer-reviewed research-aid warning requiring independent verification before citation.
\end{itemize}
\end{document}
